\documentclass{article}
\usepackage{geometry}
\usepackage{concmath}
\usepackage[OT1]{fontenc}
\usepackage{graphicx}
\usepackage{amsfonts}
\usepackage{amsmath,mathtools}
\usepackage{calrsfs}
\usepackage{color}
\usepackage[dvipsnames]{xcolor}
\usepackage[many]{tcolorbox}
\usepackage{setspace}
\usepackage{tkz-tab}               % for LINE SPACING
\usepackage{multicol}  
\geometry{margin=0.5in}
\definecolor{green}{RGB}{52, 235, 168}
\definecolor{violet}{RGB}{165, 83, 237}
\definecolor{lemon}{RGB}{170, 237, 83}
\definecolor{red}{RGB}{240, 125, 84}
\definecolor{blue}{RGB}{105, 255, 242}


\begin{document} \Large
\noindent
\colorbox{green}{Exercice 210.1}\\

\noindent
Soit \(f \in \mathcal{L}(E)\), pour \(k \geq n, f^k = 0, f^{n-1} \neq 0\)\\
ANALYSE : \\
Soit x \(\in E\) tq \((f^{n-1},...,x)\) base de E\\
Donc \(f^{n-1}(x) \neq 0, x \neq 0\)
\\
\(f^{n-1} \neq 0 \)\\
donc \(\exists x \in E\), tq \(f^{n-1}(x) \neq 0\)\\
Soit \((\lambda_0,...,\lambda_{n-1}) \in \mathbb{K}^n\) tq \begin{equation*}
    \sum_{k=0}^{n-1}\lambda_k f^k(x) = 0
\end{equation*}
\underline{Autre de formalisation}\\
par récurrence forte \(\forall k \in [|0,n-1|], \lambda_k = 0\)\\
\(\lambda_0 = 0\)\\
Soit \(k\in [|0,n-1|]\)\\
On suppose \(\lambda_0 = ... = \lambda_k = 0\)\begin{eqnarray*}
    f^{n-k-2}(\sum_{i=k+1}^{n-1} \lambda_i f^i(x)) = 0\\
    \to \lambda_{k+1} = 0
\end{eqnarray*}\\
\underline{Autre formalisation}\\
Par l'absurde : on usppose qu'il existe un \(\lambda \neq 0\)\\
On note \(k_0 min(k \text{ tq } \lambda_k \neq 0)\)\\
Donc \(\sum_{k=k_0}^ 0{n-1}\lambda_k f^k(x) =0\)\\
On compose par \(f^{n-k_0-1} \rightarrow \lambda_{k_0}f^{n-1}(x) = 0\)\\
Absurde\\

\noindent \colorbox{green}{Exercice 210.3}\\
\noindent
Soit x,y \(\in E\). (x, f(x))\\
\underline{Si (x,y) libre}\\
\(\exists \lambda_x \in \mathbb{K}, f(x) = \lambda_x x\)
\begin{align*}
    \lambda_{x+y}(x+y) &= f(x+y)\\
    &= f(x) + f(y)\\
    &= \lambda_x x + \lambda_y y
\end{align*}
donc \((\lambda_{x+y} - \lambda_x)x + (\lambda_{x+y}-\lambda_y)y = 0\)\\
donc \(\lambda_{x+y} - \lambda_x = 0\) et \(\lambda_{x+y} - \lambda_y = 0\) 
par liberté de (x,y)\\
donc \(\lambda_x = \lambda_y\)\\
\underline{Si (x,y) liée, \(x \neq 0, y \neq 0\)}
\begin{align*}
    \exists \mu \in \mathbb{K}, y = \mu x\\
    \lambda_y y &= f(y)\\
    &= f(\mu x)\\
    &= \mu f(x)\\
    &= \mu \lambda_x x 
    &= \lambda_x y
\end{align*}
Donc \(\lambda_x = \lambda_y\) car \(y \neq 0\)\\
\underline{si x = 0 et \(y \neq 0\)}\\
f(x) = 0\\

\noindent
\colorbox{green}{210.2}\\
b/ \((dim Ker f^p)_p\) est une suite croissante et majorée par dimE donc 
stationnaire.\\
On pose \(r = min({p, dim Kerf^p = Kerf^{p+1}})\)\\
Par a/, \(Kerf^r \subset Kerf^{r+1}\)\\
Et \(dim Kerf^r = dim Ker f^{r+1}\)\\
Donc \(Ker f^r = Ker f^{r+1}\)\\
c/ Mq \(\forall p \geq n, Kerf^p = Kerf^r\)\\
* Init\\
* Soit \(p\geq n\)\\
On suppose \(Kerf^p = Kerf^r\)\\
On veut montrer \(Kerf^{p+1} = Kerf^p\)\\
\(\subset\) voir a/\\
\(\supset\) Soit \(x \in Kerf^{p+1}\)\\
ie \(f^{p+1}(x) = 0\)\\
\(f^p(f(x)) = 0\)
Donc f(x) \(\in Ker^p = Kerf^r\) par HR\\
Donc \(f^{r+1}(x) = 0\)
Donc x \(\in Kerf^{r+1} = Kerf^r\)\\
d/ On a d'après le TH du rang \(dimE = dimKerf^r + Imf^r\)\\
* Montrons \(Kerf^r \cap Imf^r = {0}\)\\
Soit \(x \in Ker^r \cap Imf^r\)\\
\(f^r(x)=0\) et \(\exists t \in E\) tq \(f^r(t) = x\)\\
Donc \(f^{2r}(t) = 0\)
Donc t \(\in kerf^{2r} = Kerf^r\) d'après c/\\
Donc x = 0\\
e/ Mq \((Imf^p)_p\) est décroissante\\
Soit \(p \in \mathbb{N}\)\\
Soit x \(\in Imf^{p+1}\)
Donc \begin{align*}
    \exists t \in E, x &= f^{p+1}(t)\\
    &= f^p(f(t))\\
    &\in Imf^p
\end{align*}

\noindent \colorbox{green}{210.18}\\
Soit \((\lambda_i)_{i \in [|0,n|]}\), tq \begin{align*}
    \sum_{i=0}^{n}\lambda_i L_i &= 0\\
    \text{Soit \(k \in [|0,n|]\) } \sum_{i=0}^{n}\lambda_i L_i(\alpha_k) &= 0\\
    &= \lambda_k = 0
\end{align*}

\noindent \colorbox{green}{210.25}\\
Analyse : \\
Soit \(f \in E\) tq \(\exists f_1 \in F_1, \exists f_2 \in F_2, \exists f_3 \in F_3\)\begin{equation*}
    f = f_1 + f_2 + f_3
\end{equation*}
En particulier x = 0 \\
f(0) = \(f_1(0)\)\\
\(f_1\) constante donc \(\forall x, f(x) = f(0)\)\\
* \(x \in [-1, 0]\)\\
f(x) = f(0) + \(f_3(x)\)
etc etc 




\end{document}